\documentclass[10pt,a4paper]{amsart}
\usepackage[margin=19mm]{geometry}
\usepackage{amsmath,amssymb,mathtools}
\usepackage[scr=rsfs,cal=euler]{mathalfa}
\usepackage{xfrac}
\usepackage[unicode,hidelinks,bookmarks=false]{hyperref}
\newcommand{\ie}{i.e.\ }
\newcommand{\cf}{cf.\ }
\newcommand{\resp}{respectively\ }
\newcommand{\Subsection}{\subsection}
\providecommand{\nobreakdash}{\nobreak\hbox{-}}
\setlength{\emergencystretch}{2em}
\multlinegap0pt
\usepackage{bm}
\usepackage{bbm}
\usepackage{dsfont}
\usepackage{xspace}
\usepackage{cancel}
\usepackage{mathtools}
\usepackage{amsthm}
\makeatletter
\@ifundefined{theorem}{\newtheorem{theorem}{Theorem}}{}
\@ifundefined{lemma}{\newtheorem{lemma}[theorem]{Lemma}}{}
\@ifundefined{proposition}{\newtheorem{proposition}[theorem]{Proposition}}{}
\makeatother
\newcommand{\Int}{\operatorname{Int}}
\newcommand{\Z}{\mathbb Z}
\title[Shifted-Binomial Expansions of Dilated Binomial Polynomials:]{Shifted-Binomial Expansions of Dilated Binomial Polynomials: Multinomial Decimation, Reflection Symmetry, and Universal Divisibility}
\author{Abdelhai Doukali}
\date{Source: arXiv:2607.12173v3 (CC BY 4.0)}
\begin{document}
\maketitle

\noindent\emph{Source and licence.} Shifted-Binomial Expansions of Dilated Binomial Polynomials: Multinomial Decimation, Reflection Symmetry, and Universal Divisibility. Version arXiv:2607.12173v3 (CC BY 4.0), distributed under the Creative Commons Attribution 4.0 International licence, as recorded in the arXiv licence metadata for this exact version and shown on the abstract page https://arxiv.org/abs/2607.12173v3. Full rights notice: licenses/arxiv-2607.12173-CC-BY-4.0.txt.\par\smallskip

\noindent\textit{Editorial scope.} Counted unit: the Theorem whose source label is \texttt{thm:exact-gcd} in the article (source0.tex lines 667--707, source label \texttt{thm:exact-gcd}), delivered together with the article's own complete original proof of it (source0.tex lines 709--824, one \texttt{proof} environment of 116 lines). No mathematics is written, repaired or re-derived: statement and proof are the article's own text line for line. Required context retained uncounted: prop:support (lines 341--356); lem:fixed-divisor (lines 612--625); lem:consecutive-vp (lines 641--647). Every definition, display and label of those lines is retained unchanged. Omitted from this package and not redistributed: 1--340, 357--611, 626--640, 648--666, 708--708, 825--1133. Nothing inside a retained excerpt is omitted or altered: every retained body is delivered exactly as the article's lines are written, so no omission receipt of any kind accompanies this package. Citations: the retained excerpts carry no citation command at all, so no bibliography entry of the article is transcribed and no reference is redistributed. Reference adaptation: none was needed and none was made; every reference inside the delivered unit resolves inside it, so no unresolved reference or citation is produced. Printed slips kept literal: no printed word, symbol, hypothesis or number of the counted statement or of its proof was altered. Layout and class: the article's own class is \texttt{article} with packages including fontenc, lmodern, amsmath, amssymb, amsthm, mathtools, geometry, microtype, booktabs, enumitem, hyperref; the wrapper is the lane's pinned \texttt{amsart} editorial wrapper at 10pt on A4, which restates the theorem-like environments the retained text needs and adds the article's own macro definitions that the retained text needs (\texttt{\textbackslash Int}, \texttt{\textbackslash Z}), with extra wrapper packages (bm, bbm, dsfont, xspace, cancel, mathtools). The delivered numbering is the unit's own. Assets: no retained span contains a figure environment, a graphics-inclusion command, a PDF-page inclusion, a table or a TikZ picture; no image file is copied into this package, the package declares no asset dependency and no image file is redistributed with it. Licence: version arXiv:2607.12173v3 of this article is distributed under the Creative Commons Attribution 4.0 International licence, as recorded in the arXiv licence metadata for this exact version and shown on the abstract page https://arxiv.org/abs/2607.12173v3. A full rights notice is delivered at licenses/arxiv-2607.12173-CC-BY-4.0.txt. Characters: every retained excerpt is pure ASCII, so the delivered engine, XeLaTeX with its default Latin Modern text font, reports no missing character.
\begin{proposition}[Exact support]
\label{prop:support}
Put
\[
 K_{m,r}
 =
 d+1-\left\lceil\frac dm\right\rceil
 =
 m+r-1-\left\lceil\frac{r-1}{m}\right\rceil.
\]
Then
\[
 B_{m,r,k}>0\quad\Longleftrightarrow\quad 1\le k\le K_{m,r},
\]
and $B_{m,r,k}=0$ for $K_{m,r}<k\le d+1$.
\end{proposition}

\begin{lemma}[Fixed divisor and shifted-binomial coordinates]
\label{lem:fixed-divisor}
Let $P\in\Int(\Z)$ have degree at most $d$, and write
\[
 P(X)=\sum_{k=1}^{d+1}c_k\binom{X+k-1}{d},
 \qquad c_k\in\Z.
\]
Then
\[
 \gcd_{1\le k\le d+1}c_k
 =
 \gcd_{n\in\Z}P(n).
\]
\end{lemma}

\begin{lemma}[Minimum valuation of a product of consecutive integers]
\label{lem:consecutive-vp}
Let $p$ be prime and $L\ge1$. Then
\[
 \min_{t\in\Z}\sum_{j=0}^{L-1}v_p(t-j)=v_p(L!).
\]
\end{lemma}

\begin{theorem}[Exact row gcd]
\label{thm:exact-gcd}
Put
\[
 d=m+r-1,
 \qquad
 G_{m,r}:=\gcd_{1\le k\le K_{m,r}}B_{m,r,k}.
\]
For every prime $p\mid m$, set
\[
 a_p=v_p(m),
 \qquad
 q_p=\left\lfloor\frac{d-1}{p^{a_p}}\right\rfloor
     =\left\lceil\frac{d}{p^{a_p}}\right\rceil-1.
\]
Then
\begin{equation}
\label{eq:exact-gcd}
 G_{m,r}
 =
 \prod_{p\mid m}
 p^{\,a_p+v_p(q_p+1)-v_p(d)}.
\end{equation}
Equivalently,
\[
 G_{m,r}
 =
 \prod_{p\mid m}
 p^{
 v_p(m)
 +v_p\!\left(
 \left\lceil\frac{m+r-1}{p^{v_p(m)}}\right\rceil
 \right)
 -v_p(m+r-1)
 }.
\]
In particular,
\[
 \frac{m}{\gcd(m,r-1)}\mid G_{m,r}.
\]
\end{theorem}

\begin{proof}
By Proposition~\ref{prop:support}, the coefficients outside
$1\le k\le K_{m,r}$ vanish. Hence
\[
 G_{m,r}
 =
 \gcd_{1\le k\le d+1}B_{m,r,k}.
\]
Applying Lemma~\ref{lem:fixed-divisor} to
\[
 P(X)=\binom{mX}{d}
\]
gives
\begin{equation}
\label{eq:fixed-divisor}
 G_{m,r}
 =
 \gcd_{n\in\Z}\binom{mn}{d}.
\end{equation}
Thus it remains to compute the fixed divisor prime by prime.

Let $p$ be a prime and write
\[
 a=v_p(m).
\]
If $a=0$, multiplication by $m$ is a bijection on $\Z_p$. Therefore
\[
 \min_{n\in\Z_p}v_p\binom{mn}{d}
 =
 \min_{t\in\Z_p}v_p\binom{t}{d}.
\]
Taking $t=d$ gives value $1$, while more importantly the polynomial
$\binom{t}{d}$ has an integer value equal to $1$ at $t=d$. Hence the
minimum is $0$, and therefore
\[
 v_p(G_{m,r})=0
 \qquad(p\nmid m).
\]

Now suppose $a=v_p(m)>0$. Write
\[
 m=p^a u,\qquad p\nmid u,
\]
and put
\[
 q=\left\lfloor\frac{d-1}{p^a}\right\rfloor.
\]
For an integer $n$,
\[
 \binom{mn}{d}
 =
 \frac{\prod_{i=0}^{d-1}(mn-i)}{d!}.
\]
If $p^a\nmid i$, then $v_p(i)<a\le v_p(mn)$, and hence
\[
 v_p(mn-i)=v_p(i).
\]
For $i=p^aj$, with $0\le j\le q$,
\[
 mn-i=p^a(un-j),
\]
so
\[
 v_p(mn-i)=a+v_p(un-j).
\]
Consequently
\[
\begin{aligned}
 v_p\binom{mn}{d}
 &=
 \frac{}{}\sum_{\substack{1\le i<d\\p^a\nmid i}}v_p(i)
 +a(q+1)
 +\sum_{j=0}^{q}v_p(un-j)
 -v_p(d!).
\end{aligned}
\]
Since $u$ is coprime to $p$, multiplication by $u$ permutes the residue
classes modulo every power of $p$. Hence, for the purpose of minimizing
the finite valuation sum, $un$ may be replaced by an arbitrary integer
$t$. By Lemma~\ref{lem:consecutive-vp},
\[
 \min_n\sum_{j=0}^{q}v_p(un-j)
 =
 v_p((q+1)!).
\]
On the other hand,
\[
 v_p((d-1)!)
 =
 \sum_{\substack{1\le i<d\\p^a\nmid i}}v_p(i)
 +aq+v_p(q!).
\]
Subtracting $v_p(d!)=v_p((d-1)!)+v_p(d)$ therefore gives
\[
\begin{aligned}
 \min_n v_p\binom{mn}{d}
 &=
 a+v_p(q+1)-v_p(d).
\end{aligned}
\]
Hence
\[
 v_p(G_{m,r})
 =
 a_p+v_p(q_p+1)-v_p(d)
\]
for every $p\mid m$, while it is zero for $p\nmid m$. This proves
\eqref{eq:exact-gcd}.

Finally, write $g=\gcd(m,d)=\gcd(m,r-1)$. The exponent furnished by
\eqref{eq:exact-gcd} is at least $v_p(m)-v_p(g)$ for every $p\mid m$,
because $v_p(q_p+1)\ge0$. Therefore
\[
 \frac{m}{\gcd(m,r-1)}\mid G_{m,r}.
\]
\end{proof}

\end{document}
